\documentclass[10pt,a4paper]{amsart}
\usepackage[margin=19mm]{geometry}
\usepackage{amsmath,amssymb,mathtools}
\usepackage[scr=rsfs,cal=euler]{mathalfa}
\usepackage{xfrac}
\usepackage[unicode,hidelinks,bookmarks=false]{hyperref}
\newcommand{\ie}{i.e.\ }
\newcommand{\cf}{cf.\ }
\newcommand{\resp}{respectively\ }
\newcommand{\Subsection}{\subsection}
\providecommand{\nobreakdash}{\nobreak\hbox{-}}
\setlength{\emergencystretch}{2em}
\multlinegap0pt
\usepackage{amssymb}
\newtheorem{thm}{Theorem}
\renewcommand{\vec}[1]{\boldsymbol{#1}} % Uncomment for BOLD vectors.
\DeclareMathOperator{\wt}{wt}
\title{Second order Recurrences, quadratic number fields and cyclic codes}
\author{Minjia Shi, Xuan Wang, Bouazzaoui Zakariae, Jon-Lark Kim, Patrick Sol\'e}
\date{Source: arXiv:2603.25343v1 (CC BY 4.0)}
\begin{document}
\maketitle

\noindent\emph{Source and licence.} Second order Recurrences, quadratic number fields and cyclic codes. Version arXiv:2603.25343v1 (CC BY 4.0), distributed by arXiv under the Creative Commons Attribution 4.0 International licence, http://creativecommons.org/licenses/by/4.0/, as recorded in the arXiv licence metadata for this exact version and shown on the abstract page https://arxiv.org/abs/2603.25343v1. Full licence notice: \texttt{licenses/arxiv-2603.25343-CC-BY-4.0.txt}.\par\smallskip

\noindent\textit{Editorial scope.} Counted unit: the article's thm thm-wd (source0.tex lines 436--442). The article's complete original proof of it is delivered with it (source0.tex lines 443--479, one \texttt{proof} environment of 37 lines). No mathematics is written, repaired or re-derived: the statement and the proof are the article's own lines, in the article's own notation, and no retained line is substituted or re-flowed.  No further source-local context is needed: every cross-reference of every retained line resolves inside the two delivered spans.  Nothing was omitted from any retained span, so the package carries no omission record.  No retained line contains a citation, so the wrapper carries no bibliography and no bibliography entry.  No retained line contains a figure environment, an image-inclusion command or drawing code, so no image is delivered and the package declares no asset dependency.  Editorial formatting only, in addition: an AMS \texttt{amsart} preamble that restates the article's own declarations and macros used by the retained lines, namely the environments \texttt{thm} on separate counters, together with the standard packages those lines need.  The retained environments print in this document as Theorem 1 in order of appearance, whereas the article prints the same items with the numbering of its own full document.  The complete native source of the article as distributed in the arXiv e-print for this exact version is bundled unchanged as \texttt{sources/197200/source0.tex}.
\begin{thm} \label{thm-wd}
	Assume that $C_1$ is a two-weight MDS code whose length $n$ is at most $p+1$.
	Then
	\[ A_{n-1}(C_1) = n(p-1), \quad A_n(C_1) = (p-1)(p-n+1), \]
	and
	\[ A_{n-1}(C_2) = n(p^2-1), \quad A_n(C_2) = (p^2 - 1)(p^2 - n+1). \]
\end{thm}

\begin{proof}
	Since $C_1$ is a two-weight MDS code, then $w_1 = n-1$, $w_2 = n$, and
	\[ A_{n-1}(C_1) = n(p-1), \quad A_n(C_1) = (p-1)(p-n+1). \]
	Consider the map $\pi: C_2 \rightarrow C_1$ defined as $\pi(\vec{c}) = \vec{c} \pmod{p}$ for each $\vec{c} \in C_2$.
	So $\pi$ is homomorphism.
	It is clear that for each $\vec{c} \in C_2$, it can be uniquely expressed by $\vec{c}_1 \oplus p \vec{c}_2$, where $\vec{c}_1 = \pi(\vec{c}),\vec{c}_2 \in C_1$, and $\oplus$ means the addition operation over $\mathbb{Z}_{p^2}$.
	Obviously, the Hamming weight of $\vec{c}$ is at least $\pi(\vec{c})$, i.e., $\wt(\vec{c}) \geqslant \wt(\pi(\vec{c}))$.
	Thus, we have the following three cases:
	\begin{enumerate}
		\item[(1)] If $\pi(\vec{c}) = \vec{0}$, then $\vec{c} \in pC_2 = \langle pf_1 \rangle$ since $C_2$ is free.
		Hence, there are $n(p-1)$ such $\vec{c}$ of Hamming weight $n-1$, and $(p-1)(p-n+1)$ such $\vec{c}$ of Hamming weight $n$.
		
		\item[(2)] If $\pi(\vec{c}) \neq \vec{0}$ has Hamming weight $n$, then $\vec{c}$ has Hamming weight $n$.
		And there are $(p-1)(p-n+1)$ such $\vec{c}$ of Hamming weight $n$.
		
		\item[(3)] If $\pi(\vec{c}) \neq \vec{0}$ has Hamming weight $n-1$, then $\vec{c}$ has Hamming weight at least $n-1$.
		Note that the row of the generator matrix (in standard form) of $C_2$ has Hamming weight $n-1$.
		We claim that $\vec{c}$ has Hamming weight $n-1$ if and only if $\vec{c}_2 = \lambda \vec{c}_1 \pmod{p}$ for some $\lambda \in \mathbb{Z}_p$.
		For otherwise, we have the following two cases.
		\begin{enumerate}
			\item If $\vec{c}_2$ has Hamming weight $n$, then it is easy to check that the Hamming weight of $\vec{c}_1 \oplus p \vec{c}_2$ is exactly $n$.
			
			\item If $\vec{c}_2$ has Hamming weight $n-1$ but $\vec{c}_2 \neq \vec{c}_1$ for each nonzero $\lambda \in \mathbb{Z}_p$, then the Hamming weight of $\vec{c}_1 \oplus p \vec{c}_2$ is still exactly $n$.
		\end{enumerate}
		Hence,
		\begin{align*}
			|\{ \vec{c} \in C_2: \wt(\vec{c}) = n-1 \}| & = |\{ (\vec{c}_1, \lambda \vec{c}_1): \vec{c}_1 \in C_1, \wt(\vec{c}_1) = n-1 \}| \\
			& = np(p-1).
		\end{align*}
		the number of codewords $\vec{c}$ of Hamming weight $n-1$ is $n(p-1)$.
	\end{enumerate}
	Therefore, $C_2$ is a two-Hamming-weight code, and
	\[ A_{n-1}(C_2) = n(p-1) + np(p-1) = n(p+1)(p-1) = n(p^2 - 1), \]
	which implies that
	\[ A_{n}(C_2) = p^4 - 1 - A_{n-1} = (p^2 - 1)(p^2 + 1 - n). \]
	We complete the proof.
\end{proof}

\end{document}
